% requer \usepackage{booktabs}. Tabela descritiva.
\begin{table}[htbp]\centering
\caption{Composicao do momento da primeira solicitacao por quartil de renda do CEP e coorte de nascimento}
\label{tab:S7b}
\small
\begin{tabular}{lrrrrr}
\toprule
Coorte & Solic. & \% $<$6m & \% $<$12m & \% $<$18m & \% $<$24m \\
\midrule
\multicolumn{6}{l}{\textit{Quartil de renda do CEP 1 (mais baixa)}}}} \\
2010 & 31,235 & 6.1 & 16.3 & 31.2 & 47.0 \\
2011 & 36,029 & 15.2 & 33.6 & 49.8 & 61.2 \\
2012 & 37,581 & 27.2 & 44.4 & 58.3 & 68.3 \\
2013 & 37,298 & 30.4 & 48.6 & 62.0 & 72.0 \\
2014 & 37,200 & 33.9 & 53.2 & 68.2 & 78.3 \\
2015 & 36,231 & 38.8 & 60.3 & 74.7 & 84.6 \\
2016 & 31,722 & 44.8 & 67.2 & 82.7 & 93.2 \\
\midrule
\multicolumn{6}{l}{\textit{Quartil de renda do CEP 2}}}} \\
2010 & 31,724 & 6.1 & 15.7 & 30.7 & 46.6 \\
2011 & 36,258 & 15.5 & 33.4 & 49.7 & 61.2 \\
2012 & 37,840 & 26.8 & 44.3 & 58.2 & 68.0 \\
2013 & 37,295 & 30.1 & 47.7 & 61.3 & 71.0 \\
2014 & 37,129 & 33.2 & 52.0 & 67.1 & 77.6 \\
2015 & 35,822 & 38.6 & 59.9 & 74.5 & 84.2 \\
2016 & 31,536 & 44.2 & 66.4 & 82.0 & 92.8 \\
\midrule
\multicolumn{6}{l}{\textit{Quartil de renda do CEP 3}}}} \\
2010 & 31,921 & 5.8 & 14.7 & 29.6 & 45.9 \\
2011 & 36,384 & 15.7 & 33.3 & 49.3 & 61.2 \\
2012 & 38,614 & 27.0 & 43.4 & 57.5 & 67.6 \\
2013 & 37,795 & 29.6 & 47.3 & 60.7 & 70.8 \\
2014 & 37,120 & 32.8 & 51.2 & 65.9 & 76.6 \\
2015 & 35,585 & 38.0 & 58.6 & 73.3 & 83.5 \\
2016 & 31,321 & 43.8 & 65.5 & 81.5 & 92.5 \\
\midrule
\multicolumn{6}{l}{\textit{Quartil de renda do CEP 4 (mais alta)}}}} \\
2010 & 31,251 & 5.0 & 12.5 & 26.5 & 42.8 \\
2011 & 35,697 & 14.5 & 31.2 & 46.8 & 58.3 \\
2012 & 38,954 & 25.5 & 41.2 & 55.0 & 65.3 \\
2013 & 38,379 & 28.0 & 43.8 & 57.3 & 67.7 \\
2014 & 38,076 & 30.4 & 47.4 & 62.7 & 74.1 \\
2015 & 36,625 & 35.6 & 55.1 & 70.5 & 81.4 \\
2016 & 31,922 & 41.4 & 62.5 & 79.2 & 91.4 \\
\bottomrule
\end{tabular}
\\[3pt]{\footnotesize \textit{Note:} Quatro paineis, um por quartil de renda do CEP (Q1 mais baixa). Denominador = solicitantes da celula (coorte $\times$ quartil); numerador = os com primeira solicitacao antes da idade indicada. Ao contrario da Tabela S7a, o denominador aqui e o de solicitantes (nascidos vivos nao estao disponiveis por quartil de renda do CEP), portanto compara a composicao do momento entre grupos de renda.}
\end{table}
